% !TEX root = ../../assign.tex
\subsection{Problem 2}
\subsubsection{(\rom{1})}
\begin{definition}
    A set $S \subseteq \R$ is called inductive if $1 \in S$ and $k \in S$ implies $k + 1 \in S$. 
\end{definition}

$\R$ is inductive since $1 \in \R$ and for all $x \in \R$, $x + 1 \in \R$. \\

$\paren{0, \infty}$ is inductive since $1 \in \paren{0, \infty}$ and for all $x \in \paren{0, \infty}$, $x + 1 \in \paren{0, \infty}$. \\

$\paren{0, \infty} \setminus \set{5}$ is not inductive since $4 \in \paren{0, \infty} \setminus \set{5} \text{ but } 4 + 1 = 5 \not\in \paren{0, \infty} \setminus \set{5}$. 

\newpage
\subsubsection{(\rom{2})}
\begin{proof}
    Suppose that $S$ and $T$ are inductive sets, and $U = S \cap T$. 

    By the definition of inductive set, we have 
    \begin{equation} \label{eq1.2.1}
        1 \in S \quad \text{ and } \quad k \in S \implies k + 1 \in S. 
    \end{equation}

    and 
    \begin{equation} \label{eq1.2.2}
        1 \in T \quad \text{ and } \quad k \in T \implies k + 1 \in T. 
    \end{equation}

    Since $1$ is in both $S$ and $T$, it means that $1 \in U$. 

    By statements \eqref{eq1.2.1} and \eqref{eq1.2.2}, we get 
    \begin{equation*}
        k \in U
        \iff (k \in S \wedge k \in T) \quad
        \implies \quad (k+1 \in S \wedge k+1 \in T)
        \iff k+1 \in U.
    \end{equation*}

    Thus, $U$ is an inductive set. 
\end{proof}

\newpage
\subsubsection{(\rom{3})}
\begin{proof}
    Let $\mathcal{I }$ be the collection of all inductive subsets. 
    We define 
    \begin{equation*}
        \N = \bigcap_{S \in \mathcal{I}} S. 
    \end{equation*}

    Since 
    \begin{equation*}
        \paren{\forall S \in \mathcal{I }} 1 \in S, 
    \end{equation*}
    it follows that $1 \in \N$. 

    It suffices to show that 
    \begin{equation*}
        k \in \N \quad \implies \quad k + 1 \in \N. 
    \end{equation*}

    Suppose $k \in \N$. By the definition of the intersection of sets, it means that 
    \begin{equation*}
        (\forall S \in \mathcal{I })k \in S. 
    \end{equation*}

    It follows the definition of inductive set that 
    \begin{equation*}
        (\forall S \in \mathcal{I })k + 1 \in S, 
    \end{equation*} 

    which implies 
    \begin{equation*}
        k + 1 \in \N.  
    \end{equation*}

    Thus, $\N$ is inductive. 
\end{proof}